% GENERATED FILE. Edit canonical lessons, then run npm run generate:books.
% canonical-source-hash: fnv1a64:a9f3354716f16b1b
% canonical-lessons: JA-C482-kugi, JA-C482-nokogiri, JA-C482-harigane, JA-C482-mokuzou, JA-C482-sasaru

\chapter{Nails, Saws and Wire}
\label{ch:ja-nails-saws-and-wire}
\hlchaptermodality{482}

\begin{chapteropening}
\textbf{By the end of this chapter:} \emph{I can name the things a carpenter uses.}

Complete the last lesson of chapter 482: I can name the things a carpenter uses.
\end{chapteropening}

\section[\texorpdfstring{\ja{くぎ} (kugi)}{くぎ (kugi)}]{\ja{くぎ} (kugi) — a nail}

\label{lesson:JA-C482-kugi}



\begin{quote}
Before the new one: say the Japanese for wisdom, then the Japanese for to devise.
\end{quote}

\subsection*{You'll want to know: \ja{くぎ}}
\textbf{\ja{くぎ}} — \emph{kugi} — \textquotedblleft{}a nail\textquotedblright{}.

\begin{quote}
\textbf{\ja{いたに} \ja{くぎを} \ja{うちます。}} — \emph{ita ni kugi o uchimasu} — \textquotedblleft{}I drive a nail into the board.\textquotedblright{}
\end{quote}

\begin{cousinweb}
\textbf{\ja{くぎ}} (\emph{kugi}) is a native Japanese word (wago), inherited rather than borrowed from Chinese; its kanji was attached later as a kun'yomi, and its older origin is not settled (JMdict, KANJIDIC2).
\end{cousinweb}

\begin{grammarlens}[title={What you've built}]
One more word for describing what you lived and felt.
\end{grammarlens}

\subsection*{Your turn}
\begin{itemize}
\raggedright
  \item \emph{Say it:} \emph{kugi}
  \item \emph{Say it:} \emph{kugi}, once more
  \item \emph{Say it:} \emph{ita ni kugi o uchimasu} — I drive a nail into the board.
  \item \emph{Recall:} say the Japanese for wisdom, then the Japanese for to devise, then say \emph{kugi} again
\end{itemize}

\begin{tcolorbox}[breakable,colback=teal!4,colframe=teal!35!black,title={Before you move on}]
What does \ja{くぎ} mean? (\textquotedblleft{}A nail\textquotedblright{}.) Say it once more.
\end{tcolorbox}
\section[\texorpdfstring{\ja{のこぎり} (nokogiri)}{のこぎり (nokogiri)}]{\ja{のこぎり} (nokogiri) — a saw}

\label{lesson:JA-C482-nokogiri}



\begin{quote}
Before the new one: say the Japanese for to devise, then the Japanese for a nail.
\end{quote}

\subsection*{You'll want to know: \ja{のこぎり}}
\textbf{\ja{のこぎり}} — \emph{nokogiri} — \textquotedblleft{}a saw\textquotedblright{}.

\begin{quote}
\textbf{\ja{ちちは} \ja{のこぎりを} \ja{つかいます。}} — \emph{chichi wa nokogiri o tsukaimasu} — \textquotedblleft{}My father uses a saw.\textquotedblright{}
\end{quote}

\begin{cousinweb}
\textbf{\ja{のこぎり}} (\emph{nokogiri}) is a native Japanese word (wago), inherited rather than borrowed from Chinese, and usually written in kana; its older origin is not settled (JMdict).
\end{cousinweb}

\begin{grammarlens}[title={What you've built}]
One more word for describing what you lived and felt.
\end{grammarlens}

\subsection*{Your turn}
\begin{itemize}
\raggedright
  \item \emph{Say it:} \emph{nokogiri}
  \item \emph{Say it:} \emph{nokogiri}, once more
  \item \emph{Say it:} \emph{chichi wa nokogiri o tsukaimasu} — My father uses a saw.
  \item \emph{Recall:} say the Japanese for to devise, then the Japanese for a nail, then say \emph{nokogiri} again
\end{itemize}

\begin{tcolorbox}[breakable,colback=teal!4,colframe=teal!35!black,title={Before you move on}]
What does \ja{のこぎり} mean? (\textquotedblleft{}A saw\textquotedblright{}.) Say it once more.
\end{tcolorbox}
\section[\texorpdfstring{\ja{はりがね} (harigane)}{はりがね (harigane)}]{\ja{はりがね} (harigane) — wire}

\label{lesson:JA-C482-harigane}



\begin{quote}
Before the new one: say the Japanese for a nail, then the Japanese for a saw.
\end{quote}

\subsection*{You'll want to know: \ja{はりがね}}
\textbf{\ja{はりがね}} — \emph{harigane} — \textquotedblleft{}wire\textquotedblright{}.

\begin{quote}
\textbf{\ja{はりがねで} \ja{はこを} \ja{むすびます。}} — \emph{harigane de hako o musubimasu} — \textquotedblleft{}I tie the box up with wire.\textquotedblright{}
\end{quote}

\begin{cousinweb}
\textbf{\ja{はりがね}} (\emph{harigane}) is a native Japanese compound (wago), inherited rather than borrowed from Chinese: built from \emph{hari}, \textquotedblleft{}needle\textquotedblright{} and \emph{gane}, \textquotedblleft{}metal\textquotedblright{} (JMdict, KANJIDIC2).
\end{cousinweb}

\begin{grammarlens}[title={What you've built}]
One more word for describing what you lived and felt.
\end{grammarlens}

\subsection*{Your turn}
\begin{itemize}
\raggedright
  \item \emph{Say it:} \emph{harigane}
  \item \emph{Say it:} \emph{harigane}, once more
  \item \emph{Say it:} \emph{harigane de hako o musubimasu} — I tie the box up with wire.
  \item \emph{Recall:} say the Japanese for a nail, then the Japanese for a saw, then say \emph{harigane} again
\end{itemize}

\begin{tcolorbox}[breakable,colback=teal!4,colframe=teal!35!black,title={Before you move on}]
What does \ja{はりがね} mean? (\textquotedblleft{}Wire\textquotedblright{}.) Say it once more.
\end{tcolorbox}
\section[\texorpdfstring{\ja{もくぞう} (mokuzō)}{もくぞう (mokuzō)}]{\ja{もくぞう} (mokuzō) — made of wood; wooden}

\label{lesson:JA-C482-mokuzou}



\begin{quote}
Before the new one: say the Japanese for a saw, then the Japanese for wire.
\end{quote}

\subsection*{You'll want to know: \ja{もくぞう}}
\textbf{\ja{もくぞう}} — \emph{mokuzō} — \textquotedblleft{}made of wood; wooden\textquotedblright{}.

\begin{quote}
\textbf{\ja{この} \ja{うちは} \ja{もくぞうです。}} — \emph{kono uchi wa mokuzō desu} — \textquotedblleft{}This house is made of wood.\textquotedblright{}
\end{quote}

\begin{cousinweb}
\textbf{\ja{もくぞう}} (\emph{mokuzō}) is Sino-Japanese (kango): two Chinese characters read with their on'yomi, readings borrowed from Chinese. They are \emph{moku}, \textquotedblleft{}tree\textquotedblright{} (Middle Chinese \emph{muwk}) and \emph{zō}, \textquotedblleft{}create\textquotedblright{} (\emph{dzawX}) (KANJIDIC2; Middle Chinese in Baxter's transcription).
\end{cousinweb}

\begin{grammarlens}[title={What you've built}]
One more word for describing what you lived and felt.
\end{grammarlens}

\subsection*{Your turn}
\begin{itemize}
\raggedright
  \item \emph{Say it:} \emph{mokuzō}
  \item \emph{Say it:} \emph{mokuzō}, once more
  \item \emph{Say it:} \emph{kono uchi wa mokuzō desu} — This house is made of wood.
  \item \emph{Recall:} say the Japanese for a saw, then the Japanese for wire, then say \emph{mokuzō} again
\end{itemize}

\begin{tcolorbox}[breakable,colback=teal!4,colframe=teal!35!black,title={Before you move on}]
What does \ja{もくぞう} mean? (\textquotedblleft{}Made of wood; wooden\textquotedblright{}.) Say it once more.
\end{tcolorbox}
\section[\texorpdfstring{\ja{ささる} (sasaru)}{ささる (sasaru)}]{\ja{ささる} (sasaru) — to stick into; to pierce}

\label{lesson:JA-C482-sasaru}



\begin{quote}
Before the new one: say the Japanese for wire, then the Japanese for made of wood.
\end{quote}

\subsection*{You'll want to know: \ja{ささる}}
\textbf{\ja{ささる}} — \emph{sasaru} — \textquotedblleft{}to stick into; to pierce\textquotedblright{}.

\begin{quote}
\textbf{\ja{いたの} \ja{くぎが} \ja{てに} \ja{ささりました。}} — \emph{ita no kugi ga te ni sasarimashita} — \textquotedblleft{}A nail in the board went into my hand.\textquotedblright{}
\end{quote}

Polite \textbf{\ja{ささります}}, \emph{sasarimasu}; a godan verb. Past (polite): \textbf{\ja{くぎが} \ja{ゆびに} \ja{ささりました。}}, \emph{kugi ga yubi ni sasarimashita}.

\begin{cousinweb}
\textbf{\ja{ささる}} (\emph{sasaru}) is a native Japanese word (wago), inherited rather than borrowed from Chinese; its kanji was attached later as a kun'yomi, and its older origin is not settled (JMdict, KANJIDIC2).
\end{cousinweb}

\begin{grammarlens}[title={What you've built}]
One more word for describing what you lived and felt.
\end{grammarlens}

\subsection*{Your turn}
\begin{itemize}
\raggedright
  \item \emph{Say it:} \emph{sasaru}
  \item \emph{Say it:} \emph{sasaru}, once more
  \item \emph{Say it:} \emph{ita no kugi ga te ni sasarimashita} — A nail in the board went into my hand.
  \item \emph{Recall:} say the Japanese for wire, then the Japanese for made of wood, then say \emph{sasaru} again
\end{itemize}

\begin{tcolorbox}[breakable,colback=teal!4,colframe=teal!35!black,title={Before you move on}]
What does \ja{ささる} mean? (\textquotedblleft{}To stick into; to pierce\textquotedblright{}.) Say it once more.
\end{tcolorbox}
